\section{Каталог свидетельств: 55 работ}\label{app:catalog}

Каталог содержит сведения о 55 работах, включённых в обзор. Группа 1 — модели на понятиях; группа 2 — обучение с другими человеческими сигналами. Пометка «пограничная» означает, что связь результата с содержанием обучающего сигнала заявлена косвенно; критерии описаны в приложении C. «Не найден» ограничено указанным объёмом чтения. «Не зафиксирован в таблице извлечения» означает отсутствие сведений в записи. Полная повторная сверка всех первоисточников при этой редакции не проводилась; определения этапов HiCEM и несколько адресных формулировок уточнены. Английские цитаты сохранены.


\Needspace{9\baselineskip}\subsection*{1. \href{https://doi.org/10.48550/arxiv.2209.09056}{Concept Embedding Models: Beyond the Accuracy-Explainability Trade-Off}}

\textbf{Год и группа.} 2022; группа 1; пограничная: нет

\textbf{Обучающий сигнал A.} 28 из 112 атрибутов CUB (25\% аннотаций, случайно)

\textbf{Оценка B.} остальные 84 атрибута (линейная проба по всему узкому месту)

\textbf{Модель и данные.} CEM и Hybrid-CBM, ResNet; CUB, CelebA и др.

\textbf{Утверждение.} the bottlenecks learnt by CEMs are considerably more interpretable and can therefore serve as powerful feature extractors

\textbf{Место в тексте.} §5.3, с. 8 PDF (детали прил. A.9); аннотация: capturing meaningful semantics including and beyond their ground truth labels

\textbf{Контроль.} не найден

\textbf{Результат контроля.} Результат не зафиксирован в исходной таблице извлечения.

\textbf{Другие сравнения.} есть Hybrid-CBM и модель без понятий (на основной задаче); исходной предобученной модели в пробе нет; контрольной замены обучающего сигнала не найдено

\textbf{Объём чтения.} полный текст (аннотация, §1, §5.3, прил. A.9)


\Needspace{9\baselineskip}\subsection*{2. \href{https://doi.org/10.48550/arxiv.2602.23947}{Hierarchical Concept-based Interpretable Models}}

\textbf{Год и группа.} 2026; группа 1; пограничная: нет

\textbf{Обучающий сигнал A.} грубые понятия верхнего уровня (CUB, AwA2 и др.)

\textbf{Оценка B.} подпонятия, метки которых не использовали при обучении исходной CEM; после поиска метками найденных подпонятий обучают HiCEM и оценивают её предсказания по тестовой ROC-AUC

\textbf{Модель и данные.} HiCEM + Concept Splitting; MNIST-ADD, SHAPES, CUB, AwA2, PseudoKitchens, ImageNet

\textbf{Утверждение.} Concept Splitting discovers human-interpretable sub-concepts absent during training that can be used to train highly accurate HiCEMs

\textbf{Место в тексте.} аннотация, с. 1 PDF; табл. 1, с. 8-9

\textbf{Контроль.} не найден

\textbf{Результат контроля.} Контроль «HiCEM w/o concept splitting» сохраняет общие метки A, но не использует метки найденных подпонятий. Это проверка дополнительного обучения на подпонятиях, а не замена общих меток A. Случайные слова словаря — только в опросе людей как контроль названий

\textbf{Другие сравнения.} модель без разбиения (табл. 1, CUB 0.81 против 0.85); LF-CBM; опрос людей; случайного банка подпонятий нет; оценки на новых классах нет

\textbf{Объём чтения.} полный текст (аннотация, введение, §5, табл. 1)


\Needspace{9\baselineskip}\subsection*{3. \href{https://doi.org/10.48550/arxiv.2505.03361}{Interpretable Zero-shot Learning with Infinite Class Concepts}}

\textbf{Год и группа.} 2025; группа 1; пограничная: да

\textbf{Обучающий сигнал A.} понятия классов (человеческие и сгенерированные LLM), выбор по «переносимости и различимости»

\textbf{Оценка B.} классы, не использованные при обучении (ZSL/GZSL)

\textbf{Модель и данные.} InfZSL (I2CFormer); AWA2, CUB, SUN

\textbf{Утверждение.} our method not only improves accuracy over SOTAs across various datasets and methods, but also helps with interpretablity in ZSL

\textbf{Место в тексте.} §1 (вклады), с. 3 PDF; аннотация: ensuring that only the most transferable and discriminative concepts are selected

\textbf{Контроль.} не найден

\textbf{Результат контроля.} Абляция (табл. 2): без выбора по переносимости/различимости; случайных или бессмысленных понятий нет. Пограничный случай: включена по критерию косвенной связи результата с содержанием A (улучшение на классах, не использованных при обучении приписано выбранным по содержанию понятиям)

\textbf{Другие сравнения.} сравнение с другими методами ZSL; проверка на классах, не использованных при обучении есть по построению (ZSL)

\textbf{Объём чтения.} прочитаны отдельные разделы (аннотация, введение, §4, абляция)


\Needspace{9\baselineskip}\subsection*{4. \href{https://doi.org/10.1007/978-3-031-96625-5_18}{Interpretable Few-Shot Retinal Disease Diagnosis with Concept-Guided Prompting of Vision-Language Models}}

\textbf{Год и группа.} 2025; группа 1; пограничная: да

\textbf{Обучающий сигнал A.} понятия сетчаточных болезней (банк из GPT), предсказание понятий по снимку + токены понятий в промптах (только базовые классы)

\textbf{Оценка B.} новые классы заболеваний (novel), zero-shot mAP; 15/14 классов своего набора и 23/23 RFMiD

\textbf{Модель и данные.} CoOp/CoCoOp/ProDA/MaPLe (ViT-B/16, CLIP) + понятия; снимки глазного дна: собственный набор (29 болезней) и RFMiD

\textbf{Утверждение.} Our method not only improves retinal disease classification but also enriches few-shot and zero-shot detection (novel disease detection)

\textbf{Место в тексте.} аннотация, с. 1 PDF; табл. 2, §4.2 (+2.7 п. mAP в среднем); введение: «enhances interpretability, even for previously unseen retinal disease categories»

\textbf{Контроль.} не найден

\textbf{Результат контроля.} Результат не зафиксирован в исходной таблице извлечения.

\textbf{Другие сравнения.} те же методы без понятий (CoOp, CoCoOp, ProDA, MaPLe), Image-only ViT и CLIP linear probe; абляции числа и позиции токенов понятий; контрольных меток понятий (случайных или перемешанных) не найдено. Пограничный случай: улучшение приписано включению знаний о понятиях (банк GPT охватывает все болезни, включая новые — табл. 2 текст: «entire concept bank … was included during training»), то есть текстовые описания включали и понятия, относящиеся к оценке B; косвенная связь результата с содержанием A

\textbf{Объём чтения.} полный текст (аннотация, введение, §4.2, табл. 2, абляции), arXiv 2503.02917


\Needspace{9\baselineskip}\subsection*{5. \href{https://doi.org/10.18653/2026.emnlp-main.1709}{Learning Concepts, Not Tokens: Self-Supervised Semantic Alignment for Language Models}}

\textbf{Год и группа.} 2026; группа 2; пограничная: нет

\textbf{Обучающий сигнал A.} понятия как множества контекстных синонимов; цель предсказания понятий вместе с NTP (λ — вес)

\textbf{Оценка B.} сходство суждений людей (STS, 9 наборов), задачи MTEB (классификация, кластеризация, ранжирование), рассуждения, перплексия

\textbf{Модель и данные.} языковые модели \textasciitilde{}1B–8B (Qwen, Falcon и др.); дообучение на C4 / OpenWebText, 3 запуска; arXiv 2603.29123

\textbf{Утверждение.} This contrast supports our claim that the gains come from semantically meaningful concept sets, not from adding noise.

\textbf{Место в тексте.} §4.1, с. 5 PDF; рис. 3, 16; аннотация, с. 1

\textbf{Контроль.} есть

\textbf{Результат контроля.} «Randomized synonyms» — каждое истинное множество синонимов заменено случайным того же размера из общего пула, та же функция потерь. STS при малом λ=0.25 даёт максимум, при росте λ падает (рис. 16), у настоящих — растёт. При больших λ качество с контрольными метками ниже, чем с настоящими; при малом λ часть прироста может быть регуляризацией (авторы так и пишут)

\textbf{Другие сравнения.} NTP-only FT (λ=0), предобученная модель, NTP с аугментацией синонимами (равное число шагов / эпох)

\textbf{Объём чтения.} полный текст (аннотация, §3–5, подписи рисунков, прил. D)


\Needspace{9\baselineskip}\subsection*{6. \href{https://doi.org/10.48448/tzyh-b804}{Flexible Concept Bottleneck Model}}

\textbf{Год и группа.} 2026; группа 1; пограничная: да

\textbf{Обучающий сигнал A.} заданный набор понятий (LLM / VLM-понятия) при обучении гиперсети и разрежённого выбора

\textbf{Оценка B.} новые наборы понятий, подставленные без переобучения (zero-shot) и после 1 эпохи дообучения

\textbf{Модель и данные.} FCBM: гиперсеть по эмбеддингам понятий, sparsemax; 5 наборов (в т.ч. Places365, CUB); ResNet50, ViT-L/14; arXiv 2511.06678. Пограничный случай: улучшение приписано архитектуре и «структуре семантического пространства», не содержанию A

\textbf{Утверждение.} suggesting the learned structure within FCBM captures generalizable patterns in the semantic space

\textbf{Место в тексте.} §4.3, с. 5 PDF; аннотация: generalizes well to unseen concepts

\textbf{Контроль.} не найден

\textbf{Результат контроля.} Результат не зафиксирован в исходной таблице извлечения.

\textbf{Другие сравнения.} абляции sparsemax / температуры (табл. 2); сравнение с CBM-базами; случайных или бессмысленных понятий нет

\textbf{Объём чтения.} полный текст (аннотация, §3–5, поиск ключевых слов); опубликованная версия — arXiv 2511.06678


\Needspace{9\baselineskip}\subsection*{7. \href{https://doi.org/10.48550/arxiv.2602.00924}{Supervised sparse auto-encoders for interpretable and compositional representations}}

\textbf{Год и группа.} 2026; группа 1; пограничная: да

\textbf{Обучающий сигнал A.} метки понятий (контролируемый разреженный автокодировщик: у каждого понятия свой блок, учится только на запросах с этим понятием); Stable Diffusion 3.5, вложения запросов (T5), 1500 запросов

\textbf{Оценка B.} неувиденные в обучении сочетания известных понятий (композиционное обобщение); правка на уровне признаков

\textbf{Модель и данные.} декодерный SSAE на вложениях запросов Stable Diffusion 3.5; сочетания понятий (цвет глаз, волос, поза, предметы)

\textbf{Утверждение.} The model generalizes to unseen combinations, supporting compositional modular edits.

\textbf{Место в тексте.} §4 (Experiments), с. 6 PDF; аннотация, с. 1: successfully reconstructing images with concept combinations unseen during training. Пограничный случай: B — новые сочетания тех же понятий, не новые понятия; причина приписана структуре надзора

\textbf{Контроль.} не найден

\textbf{Результат контроля.} Результат не зафиксирован в исходной таблице извлечения.

\textbf{Другие сравнения.} качественное сравнение с неконтролируемыми SAE (рис. 3); количественных контролей нет [не сверено: прил.]

\textbf{Объём чтения.} полный текст (аннотация, введение, §3–5, выводы)


\Needspace{9\baselineskip}\subsection*{8. \href{https://doi.org/10.48550/arxiv.2603.10084}{Digging Deeper: Learning Multi-Level Concept Hierarchies}}

\textbf{Год и группа.} 2026; группа 1; пограничная: нет

\textbf{Обучающий сигнал A.} грубые понятия верхнего уровня (CUB, AwA2, SHAPES, MNIST-ADD, PseudoKitchens-2)

\textbf{Оценка B.} подпонятия нескольких уровней, метки которых не использовали при обучении исходной CEM; затем на найденных метках обучают Deep-HiCEM и оценивают её тестовую ROC-AUC по сопоставленным понятиям банка. Порог 0.7 применяется при сопоставлении

\textbf{Модель и данные.} Deep-HiCEM + Multi-Level Concept Splitting; тот же протокол, что в HiCEM (Hill 2026)

\textbf{Утверждение.} MLCS discovers human-interpretable concepts absent during training

\textbf{Место в тексте.} аннотация, с. 1 PDF; табл. 1, с. 6

\textbf{Контроль.} не найден

\textbf{Результат контроля.} Результат не зафиксирован в исходной таблице извлечения.

\textbf{Другие сравнения.} сравнение с HiCEM + Concept Splitting (0.84 против 0.85 на CUB) и LF-CBM (0.77); случайного банка подпонятий и контрольной замены обучающего сигнала не найдено

\textbf{Объём чтения.} полный текст (аннотация, введение, §3–5, табл. 1–3)


\Needspace{9\baselineskip}\subsection*{9. \href{https://doi.org/10.1109/cvpr52734.2025.02820}{Attribute-formed Class-specific Concept Space: Endowing Language Bottleneck Model with Better Interpretability and Scalability}}

\textbf{Год и группа.} 2025; группа 1; пограничная: да

\textbf{Обучающий сигнал A.} языковые понятия — описания атрибутов классов (генерирует LLM), учится классификатор понятий и визуальные промпты атрибутов на базовых классах

\textbf{Оценка B.} классы, не использованные при обучении (zero-shot, base-to-novel, few-shot)

\textbf{Модель и данные.} ALBM + VAPL на CLIP; 9 наборов (Food-101, CUB, Aircraft, Flowers102 и др.)

\textbf{Утверждение.} the improvement on novel classes indicates that our attribute-formed paradigm effectively facilitates cross-category transfer

\textbf{Место в тексте.} §4.4, Ablation study, с. 7 PDF; табл. 5. Пограничный случай (как InfZSL): приписано единому набору атрибутов, а не содержанию A

\textbf{Контроль.} не найден

\textbf{Результат контроля.} Результат не зафиксирован в исходной таблице извлечения.

\textbf{Другие сравнения.} абляция (без Lw = zero-shot; +2.3 п. на новых классах); LaBo и «случайно инициализированный слой понятий» (рис. A3, с. 16) — это случайная инициализация слоя, а не обучение на контрольных метках; контрольной замены обучающего сигнала не найдено

\textbf{Объём чтения.} полный текст (аннотация, введение, §3, §4.4, прил. A–B, grep)


\Needspace{9\baselineskip}\subsection*{10. \href{https://doi.org/10.1109/ccwc51732.2021.9376108}{Explainable Neural Network-based Modulation Classification via Concept Bottleneck Models}}

\textbf{Год и группа.} 2021; группа 1; пограничная: да

\textbf{Обучающий сигнал A.} промежуточные понятия сигналов (амплитуда, частота, порядок модуляции — регрессионные головы), метки понятий известны

\textbf{Оценка B.} сигналы типов модуляции, не виденных в обучении (out-of-set): точность регрессии понятий; сам zero-shot классификатор НЕ реализован («left for future work»)

\textbf{Модель и данные.} Concept Bottleneck: CNN + головы регрессии + классификатор; синтетические данные IQ (in-set, near-set, out-of-set)

\textbf{Утверждение.} shows potential for classifying modulation schemes not seen during training (i.e. zero-shot learning)

\textbf{Место в тексте.} аннотация, с. 1 PDF; §V-C и заключение (VI)

\textbf{Контроль.} не найден

\textbf{Результат контроля.} Результат не зафиксирован в исходной таблице извлечения.

\textbf{Другие сравнения.} сравнение с одно-сетевой CNN на in-set/near-set; Пограничный случай: утверждение мягкое («promise», «likely feasible»), классификация новых типов не проверена, мерится качество предсказания понятий на сигналах, не использованных при обучении

\textbf{Объём чтения.} полный текст (аннотация, введение, §III–VI)


\Needspace{9\baselineskip}\subsection*{11. \href{https://doi.org/10.18653/v1/2024.findings-naacl.131}{PEEB: Part-based Image Classifiers with an Explainable and Editable Language Bottleneck}}

\textbf{Год и группа.} 2024; группа 1; пограничная: да

\textbf{Обучающий сигнал A.} текстовые описания частей объекта (дескрипторы птиц, GPT-4/люди), контрастное предобучение на уровне частей на Bird-11K

\textbf{Оценка B.} классы птиц и собак, не использованные при обучении (GZSL/ZSL): CUB, NABirds, iNaturalist, Dogs-120

\textbf{Модель и данные.} PEEB (OWL-ViT + CLIP-кодировщик текста, Part MLP); Bird-11K, CUB-200, NABirds, iNaturalist, Dogs-120

\textbf{Утверждение.} the expressive part descriptors play a major role in PEEB whose accuracy decreases significantly to near random-chance

\textbf{Место в тексте.} §5.1, табл. 1 (перемешанные дескрипторы: 64.33 $\rightarrow$ 0.88); §5.2–5.4 (GZSL, ZSL), прил. E.5 [номер страницы PDF не сверен]

\textbf{Контроль.} частичный

\textbf{Результат контроля.} Дескрипторы перемешаны между классами ТОЛЬКО на тесте у обученной модели (табл. 1): точность падает 64.33 $\rightarrow$ 0.88. Это проверка зависимости от смысла дескрипторов при предсказании, не контрольная замена меток при обучении; модель с контрольными дескрипторами при обучении не обучалась. Пограничный случай

\textbf{Другие сравнения.} CLIP и Menon-Vondrick (M\&V) с перемешиванием/без имён классов (табл. 1–2); необученный PEEB $\approx$ случайно (прил. E.5, табл. A8); сравнение с CBM, ProtoPNet и ZSL-методами (табл. 4–5)

\textbf{Объём чтения.} полный текст (аннотация, введение, §4–5, §7, прил. E.5)


\Needspace{9\baselineskip}\subsection*{12. \href{https://doi.org/10.1038/s41467-026-71267-5}{Human-like cognitive generalization for large models via mental representation-guided supervision}}

\textbf{Год и группа.} 2026; группа 2; пограничная: нет

\textbf{Обучающий сигнал A.} сигналы мозга (fMRI) при просмотре 150 конкретных категорий; выравнивание структуры представлений DNN с мозгом (графовое сопоставление)

\textbf{Оценка B.} 50 конкретных категорий, не использованных при обучении, и абстрактные понятия; сходство по WordNet и суждениям людей; few/zero-shot, OOD-распознавание

\textbf{Модель и данные.} CLIP-base/large, SimCLR, DINOv2 + fMRI-кодировщик; данные fMRI и изображений (150+50 категорий)

\textbf{Утверждение.} brain-in-the-loop supervised learning, utilizing a small set of brain signals, can effectively transfer human conceptual structures to DNNs

\textbf{Место в тексте.} аннотация, с. 1 PDF (arXiv 2505.09085); §Results «alignment was well generated to unseen», с. 6; обсуждение с. 13

\textbf{Контроль.} не найден

\textbf{Результат контроля.} Результат не зафиксирован в исходной таблице извлечения.

\textbf{Другие сравнения.} исходные модели без обучения на данных активности мозга (сравнение до/после, парные t-тесты); разные архитектуры и испытуемые; контрольной замены (перемешанных данных мозга или случайного графа) не найдено

\textbf{Объём чтения.} полный текст препринта arXiv 2505.09085 (версия под тем же названием «Brain-in-the-loop»; аннотация, введение, результаты, обсуждение; поиск по shuffl/random/control); слой исправлен на 2 при сводке 10.11 (A — данные fMRI)


\Needspace{9\baselineskip}\subsection*{13. \href{https://doi.org/10.1038/s41551-026-01741-4}{CLEAR: an auditable foundation model for radiology grounded in clinical concepts}}

\textbf{Год и группа.} 2026; группа 1; пограничная: да

\textbf{Обучающий сигнал A.} пары снимок–отчёт (0.87 млн, 239 391 пациент); из отчётов LLM извлекает радиологические наблюдения (понятия), снимок проецируется в пространство эмбеддингов LLM

\textbf{Оценка B.} zero-shot обнаружение патологий на четырёх внешних наборах (VinDr-CXR, PadChest, CheXpert, Indiana); линейные пробы

\textbf{Модель и данные.} CLEAR: визуальный кодировщик + проекция через эмбеддинги LLM (SFR-Embedding-Mistral и др.); обучение на MIMIC-CXR, CheXpert-Plus, ReXGradient

\textbf{Утверждение.} takes advantage of the LLM's ability to capture nuanced semantic relationships within complex medical text

\textbf{Место в тексте.} Introduction/Results, с. 3 PDF; выводы с. 8 (strong zero-shot recognition)

\textbf{Контроль.} не найден

\textbf{Результат контроля.} Результат не зафиксирован в исходной таблице извлечения.

\textbf{Другие сравнения.} сравнение с CheXzero, BiomedCLIP, OpenAI CLIP; выбор текстового эмбеддера (разные LLM); контрольной замены (случайных или перемешанных понятий либо отчётов) не найдено. Пограничный случай: включена по критерию косвенной связи результата с содержанием A — улучшение на B приписано смысловому пространству LLM, а не прямо содержанию A

\textbf{Объём чтения.} полный текст (аннотация, введение, результаты, методы; поиск по shuffl/random/control/ablat)


\Needspace{9\baselineskip}\subsection*{14. \href{https://doi.org/10.48550/arxiv.2112.10871}{Translational Concept Embedding for Generalized Compositional Zero-shot Learning}}

\textbf{Год и группа.} 2021; группа 1; пограничная: да

\textbf{Обучающий сигнал A.} отношения между понятиями (атрибут–объект) в пространстве слов: потеря RVC выравнивает расстояния эмбеддингов понятий с предобученными словесными эмбеддингами; обучение на виденных композициях

\textbf{Оценка B.} сочетания «атрибут, объект», не использованные при обучении (generalized CZSL)

\textbf{Модель и данные.} Translational Concept Embedding; ResNet-18 признаки; MIT-States, UT-Zappos и др. [не сверено по наборам]

\textbf{Утверждение.} the difference of two concepts in word embedding space is a good indicator of the relation of the two concepts in the image feature space

\textbf{Место в тексте.} §4.2 Ablation Study, с. 9 PDF; рис. 2, с. 10

\textbf{Контроль.} не найден

\textbf{Результат контроля.} Результат не зафиксирован в исходной таблице извлечения.

\textbf{Другие сравнения.} абляция членов потери (табл. 3), в т. ч. без Lrvc; варианты числа пар и margin; случайных/перемешанных словесных эмбеддингов нет. Пограничный случай: включена по критерию косвенной связи результата с содержанием A (прирост на сочетаниях, не использованных при обучении приписан знанию из словесных эмбеддингов)

\textbf{Объём чтения.} полный текст (аннотация, метод, §4.2, выводы)


\Needspace{9\baselineskip}\subsection*{15. \href{https://doi.org/10.48550/arxiv.2610.22638}{Inductive Learning of Concept Embeddings from Cross-Linguistic Colexification Patterns with Textual Attributes}}

\textbf{Год и группа.} 2026; группа 1; пограничная: да

\textbf{Обучающий сигнал A.} паттерны колексификации (граф по CLICS/CLIPS) + краткие текстовые описания понятий (Sentence-BERT); обучение индуктивных эмбеддингов понятий

\textbf{Оценка B.} понятия, не вошедшие в обучение; оценка по LSIM, SemShift, WA и категориям MultiPic

\textbf{Модель и данные.} SkipGram по случайным блужданиям + текстовый кодировщик (Sentence-BERT); CLICS$^4$/CLIPS, Concepticon

\textbf{Утверждение.} the model learns meaningful representations for the concepts seen in the training data and generalizes well to unseen concepts

\textbf{Место в тексте.} §5 Conclusion and Outlook, с. 8 PDF; рис. 3-4 (с. 7-8); табл. 3

\textbf{Контроль.} не найден

\textbf{Результат контроля.} Результат не зафиксирован в исходной таблице извлечения.

\textbf{Другие сравнения.} сравнение с трансдуктивными эмбеддингами (лучше на виденных, сопоставимы на понятиях, не использованных при обучении); контрольной замены (перемешанных колексификаций или описаний) не найдено. Пограничный случай: слой 2 (лингвистический сигнал); «обобщение» приписано индуктивному построению, смысловое объяснение мягкое

\textbf{Объём чтения.} полный текст (аннотация, введение, методы, результаты, выводы)


\Needspace{9\baselineskip}\subsection*{16. \href{https://doi.org/10.1145/3767308.3836354}{AlignFace: Human-Aligned Face Similarity Metric with Interpretable Concept Relations}}

\textbf{Год и группа.} 2026; группа 1; пограничная: да

\textbf{Обучающий сигнал A.} человеческие суждения о сходстве лиц (тройки, 2AFC, датасет FACETS) с атрибутными понятиями (CBM + нейронный GAM)

\textbf{Оценка B.} другие наборы лиц (CMU Multi-PIE, CelebA — кросс-наборная проверка), отложенные тройки

\textbf{Модель и данные.} AlignFace: VLM (CLIP/FLIP) + gated cross-attention + CBM + GAM; FACETS

\textbf{Утверждение.} perhaps due to the VLM's open-domain knowledge, the use of cognitively-grounded attributes, and the alignment of relationships

\textbf{Место в тексте.} §5.3 Cross-Dataset Generalizability, с. 8 PDF; табл. 2

\textbf{Контроль.} не найден

\textbf{Результат контроля.} Результат не зафиксирован в исходной таблице извлечения.

\textbf{Другие сравнения.} метрики LPIPS, SSIM, DISTS и др.; абляции (§5.4, рис. 10: без gating, без условия группы); внутри- против меж-наборной оценки; перемешанных суждений людей или случайных атрибутов нет. Пограничный случай: слой 2; приписывание мягкое («perhaps due to»)

\textbf{Объём чтения.} полный текст (arXiv 2608.14130: аннотация, метод, §5, обсуждение)


\Needspace{9\baselineskip}\subsection*{17. \href{https://doi.org/10.52202/075280-2208}{DreamSim: Learning New Dimensions of Human Visual Similarity using Synthetic Data}}

\textbf{Год и группа.} 2023; группа 2; пограничная: нет

\textbf{Обучающий сигнал A.} человеческие суждения «какое из двух изображений больше похоже на эталон» (триплеты 2AFC) на синтетических изображениях Stable Diffusion

\textbf{Оценка B.} реальные изображения: поиск (ImageNet-R, COCO), реконструкция (инверсия признаков), BAPPS/THINGS/IQA

\textbf{Модель и данные.} DINO+CLIP+OpenCLIP, дообучение LoRA/MLP на NIGHTS (\textasciitilde{}20k триплетов); оценка на BAPPS, THINGS, ImageNet-R, COCO

\textbf{Утверждение.} our metric generalizes to real images, giving strong results on retrieval and reconstruction tasks

\textbf{Место в тексте.} аннотация, с. 1 PDF; §5.1, с. 7 (обучение на NIGHTS улучшает BAPPS, ухудшает THINGS); §6, с. 8

\textbf{Контроль.} не найден

\textbf{Результат контроля.} Перемешанных или случайных суждений нет.

\textbf{Другие сравнения.} предобученные базовые модели; ансамбль без дообучения (§6, с. 8); MLP против LoRA; LPIPS/DISTS

\textbf{Объём чтения.} полный текст (аннотация, §1, §4–§6, прил. частично)


\Needspace{9\baselineskip}\subsection*{18. \href{https://doi.org/10.18653/v1/2020.acl-main.467}{Intermediate-Task Transfer Learning with Pretrained Language Models: When and Why Does It Work?}}

\textbf{Год и группа.} 2020; группа 2; пограничная: нет

\textbf{Обучающий сигнал A.} промежуточная задача NLP (11 задач, напр. CosmosQA, HellaSwag) — дообучение RoBERTa

\textbf{Оценка B.} 10 целевых задач и 25 проб-задач (языковые навыки), на которые модель не обучалась

\textbf{Модель и данные.} RoBERTa-large; 110 пар промежуточная–целевая; GLUE/SuperGLUE и др.

\textbf{Утверждение.} intermediate tasks requiring high-level inference and reasoning abilities tend to work best

\textbf{Место в тексте.} аннотация, с. 1 PDF; §4.2 (корреляции проб и целевых задач), с. 6–7; §6, с. 9. Мягкая форма: приписано «навыкам», которым учит промежуточная задача

\textbf{Контроль.} не найден

\textbf{Результат контроля.} Промежуточной задачи со случайными или перемешанными метками нет; в статье такого контроля не найдено.

\textbf{Другие сравнения.} базовая модель без промежуточного обучения; 25 проб-задач; авторы сами отмечают, что связь проб и целевых задач слабая (забывание)

\textbf{Объём чтения.} полный текст (аннотация, §1–§6)


\Needspace{9\baselineskip}\subsection*{19. \href{https://doi.org/10.1111/cogs.12670}{Evaluating (and Improving) the Correspondence Between Deep Neural Networks and Human Representations}}

\textbf{Год и группа.} 2018; группа 2; пограничная: нет

\textbf{Обучающий сигнал A.} человеческие оценки сходства пар изображений (6 доменов, \textasciitilde{}428 тыс. оценок); линейное преобразование признаков DNN подогнано под них

\textbf{Оценка B.} (1) другие домены/отложенные изображения (межобластной перенос, отложенные домены); (2) сложность выучивания категорий из новых изображений людьми

\textbf{Модель и данные.} AlexNet, VGG, GoogLeNet, ResNet (признаки ImageNet); линейное преобразование с регуляризацией; MTurk

\textbf{Утверждение.} The resulting models transfer to new stimuli, and can be used to predict complex behaviors such as the ease of category learning.

\textbf{Место в тексте.} Discussion, с. 15 PDF; §3.5 (перенос между доменами, с. 12); §4 (опыт 2, с. 12–14)

\textbf{Контроль.} частичный

\textbf{Результат контроля.} §3.4 (с. 11–12): базовые модели, где строки/столбцы матрицы признаков случайно перемешаны при тех же оценках сходства и той же кросс-проверке: R² < 0.01 (контроль переобучения на самих A, не проверка на B). Для переноса на B (оценка по домену, категории) контрольной замены обучающего сигнала не найдено.

\textbf{Другие сравнения.} необработанные признаки DNN; HOG+SIFT; контрольная кросс-проверка без общих изображений («CV Control», табл. 1); перенос между доменами хуже, чем у необработанных признаков

\textbf{Объём чтения.} полный текст (аннотация, §1–§5)


\Needspace{9\baselineskip}\subsection*{20. \href{https://doi.org/10.1038/s41586-025-09631-6}{Aligning machine and human visual representations across abstraction levels}}

\textbf{Год и группа.} 2025; группа 2; пограничная: да

\textbf{Обучающий сигнал A.} человеческие суждения «лишнее из трёх» (THINGS) $\rightarrow$ учитель-модель $\rightarrow$ синтетические триплеты AligNet (миллионы) для дообучения студентов

\textbf{Оценка B.} новые наборы суждений людей на разных уровнях абстракции (Levels), few-shot, устойчивость к сдвигу распределения (BREEDS и др.)

\textbf{Модель и данные.} ViT, SigLIP, DINOv2 и др. (дообучение); текст прочитан по arXiv-версии 2409.06509v4 (в Nature — та же работа)

\textbf{Утверждение.} They also perform better on a diverse set of machine learning tasks, increasing generalization and out-of-distribution robustness.

\textbf{Место в тексте.} аннотация, с. 1 PDF; §2.3 «Alignment improves generalization and out-of-distribution robustness», с. 9; прил. B.2.7 (абляции), с. 76–78 PDF

\textbf{Контроль.} частичный

\textbf{Результат контроля.} Пограничный случай. Нет перемешанных человеческих суждений. Есть (1) UnAligNet — тот же конвейер и та же функция потерь, но учитель без выравнивания по человеческим данным: прирост few-shot в основном остаётся (разница мала), выравнивание даёт стабильное небольшое преимущество, ANOVA p < 0.01 (с. 76); (2) случайная выборка триплетов («as a control», с. 78) — заметно хуже кластерной. Оба контроля не есть «бессмысленные метки» в строгом смысле.

\textbf{Другие сравнения.} не дообученная базовая модель; случайно инициализированный ViT-B («Scratch»); класс-граничная выборка триплетов; размер модели

\textbf{Объём чтения.} arXiv-версия: аннотация, §1–§2, прил. B.2 (абляции); Nature-версию не открывал


\Needspace{9\baselineskip}\subsection*{21. \href{https://doi.org/10.48550/arxiv.2211.01201}{Human alignment of neural network representations}}

\textbf{Год и группа.} 2022; группа 2; пограничная: нет

\textbf{Обучающий сигнал A.} человеческие суждения «лишнее из трёх» (THINGS, Hebart 2020): линейное преобразование представлений подогнано под них

\textbf{Оценка B.} выравнивание с другими наборами человеческого сходства (Cichy 2019, King 2019) и отложенные изображения THINGS

\textbf{Модель и данные.} \textasciitilde{}100 предобученных моделей (ImageNet, CLIP и др.); линейный зонд; оценка odd-one-out и RSA

\textbf{Утверждение.} Linear transformations of neural network representations learned from behavioral responses from one dataset substantially improve alignment with human similarity judgments on the other two datasets.

\textbf{Место в тексте.} аннотация, с. 1 PDF; §4.3, с. 6 (зонд обучен только на THINGS)

\textbf{Контроль.} не найден

\textbf{Результат контроля.} Преобразования на перемешанных суждениях не проверялись; в статье такого контроля не найдено.

\textbf{Другие сравнения.} сырые (непреобразованные) представления; сравнение по обучающим данным, целям и размеру модели

\textbf{Объём чтения.} полный текст (аннотация, §1–§4, часть приложений). Утверждение мягкое: B — тот же тип человеческого сигнала; «понятия человека» упомянуты, но приписывание смыслу A слабое


\Needspace{9\baselineskip}\subsection*{22. \href{https://doi.org/10.1162/tacl_a_00363}{Infusing Finetuning with Semantic Dependencies}}

\textbf{Год и группа.} 2021; группа 2; пограничная: да

\textbf{Обучающий сигнал A.} лингвистическая разметка: семантические зависимости (предикатно-аргументная структура, DM) через парсер и RGCN в дообучении (SIFT)

\textbf{Оценка B.} задачи GLUE (8 задач; MNLI, CoLA, RTE и др.), диагностика HANS

\textbf{Модель и данные.} RoBERTa-base/large + RGCN поверх семантических графов; парсер обучен на CoNLL 2019

\textbf{Утверждение.} These results provide evidence supporting our hypothesis that incorporating semantic structures is more beneficial to RoBERTa than syntactic ones.

\textbf{Место в тексте.} §5 (результаты GLUE, табл. 3), с. 7 PDF; аннотация, с. 1; §8, с. 10

\textbf{Контроль.} не найден

\textbf{Результат контроля.} Случайных или перемешанных графов нет. Сравнение с графами синтаксических зависимостей — другой осмысленный обучающий сигнал. Пограничный случай: A служит входом дообучения, а не целью обучения; B — GLUE.

\textbf{Другие сравнения.} RoBERTa без графов; SIFT-Light без парсинга при выводе; модель со синтаксическим графом; GCN/GAT/Hidden/Scaffold (§6, прил. D)

\textbf{Объём чтения.} аннотация, §1, §4–§6, §8, прил. D (по ключевым словам)


\Needspace{9\baselineskip}\subsection*{23. \href{https://doi.org/10.52202/075280-2218}{Improving neural network representations using human similarity judgments}}

\textbf{Год и группа.} 2023; группа 2; пограничная: нет

\textbf{Обучающий сигнал A.} человеческие суждения «лишнее из трёх» (THINGS) $\rightarrow$ линейное преобразование gLocal (глобальное выравнивание + локальная контрастная потеря)

\textbf{Оценка B.} few-shot классификация (CIFAR-100, SUN397, BREEDS и др.), обнаружение аномалий — задачи, которых в обучении преобразования не было

\textbf{Модель и данные.} предобученные ViT/ResNet и image/text модели (CLIP и др.); текст прочитан по arXiv 2306.04507v2 (NeurIPS 2023)

\textbf{Утверждение.} incorporating this global structure into neural network representations improves performance on downstream tasks

\textbf{Место в тексте.} аннотация, с. 1 PDF; §4.3 few-shot, с. 7; §4.4 аномалии; табл. 1–3

\textbf{Контроль.} есть

\textbf{Результат контроля.} тройки «наперекор людям» (human-adversarial: лишним выбран другой объект, чем у людей) — пространство и качество не меняются или заметно хуже (§5 Discussion, прил. D.4); исправлено вручную 10.11 по первоисточнику (arXiv 2306.04507), агент партии 8 это место пропустил

\textbf{Другие сравнения.} наивное преобразование (только глобальная потеря) ухудшает результат; непреобразованные представления; сочетание с методами для few-shot (прил., табл. 7)

\textbf{Объём чтения.} аннотация, §1, §3–§4, прил. (поиск по ключевым словам); §5 и прил. D.4 сверены вручную 10.11


\Needspace{9\baselineskip}\subsection*{24. \href{https://doi.org/10.52202/079017-1756}{When does perceptual alignment benefit vision representations?}}

\textbf{Год и группа.} 2024; группа 2; пограничная: да

\textbf{Обучающий сигнал A.} суждения людей о сходстве: 13 900 троек изображений NIGHTS (выбор «лишнего»)

\textbf{Оценка B.} другие задачи: счёт, сегментация, оценка глубины, поиск экземпляров, RAG, VTAB, мед. снимки

\textbf{Модель и данные.} CLIP, DINO, DINOv2, SynCLR, LoRA-дообучение; NIGHTS; проверка на стандартных бенчмарках

\textbf{Утверждение.} even on tasks requiring skills beyond those needed during training or perceptual alignment

\textbf{Место в тексте.} §1, с. 2 PDF; аннотация, с. 1: injecting an inductive bias about human perceptual knowledge into vision models can contribute to better representations

\textbf{Контроль.} частичный

\textbf{Результат контроля.} Пограничный случай. §4.5, с. 8 PDF, прил. A.2, C.4: те же размер и потеря, но тройки BAPPS (искажения), THINGS (понятия) и «ImageNet» (две картинки одного класса против одной другого) — не перемешанные суждения, а другие осмысленные сигналы. «ImageNet»-тройки по классу почти ничего не дают, BAPPS минимальный эффект, THINGS ухудшает. Тройки со случайными метками нет

\textbf{Другие сравнения.} исходные предобученные модели; число шагов дообучения; SynCLR (предобучена на синтетике) как проверка, что падение — от выравнивания; доля данных

\textbf{Объём чтения.} полный текст (аннотация, §1, §4.5, §5, прил. A.2, C.4)


\Needspace{9\baselineskip}\subsection*{25. \href{https://doi.org/10.1162/imag_a_00525}{Alignment of auditory artificial networks with massive individual fMRI brain data leads to generalisable improvements in brain encoding and downstream tasks}}

\textbf{Год и группа.} 2025; группа 2; пограничная: нет

\textbf{Обучающий сигнал A.} данные fMRI 6 человек, смотревших 3 сезона «Друзей» (дообучение SoundNet, чтобы предсказывать измеренную активность мозга)

\textbf{Оценка B.} 4-й сезон (мозг) и 19 задач бенчмарка HEAR (звуковые задачи с отдельным обучением пробы)

\textbf{Модель и данные.} SoundNet (\textasciitilde{}2,5 млн параметров); Courtois NeuroMod, HEAR

\textbf{Утверждение.} brain aligning a pretrained CNN network led to more generalisable representations

\textbf{Место в тексте.} Discussion, подраздел о HEAR (версия Imaging Neuroscience, PMC12319826, XML без нумерации страниц); аннотация: brain alignment ... generalisable representations

\textbf{Контроль.} не найден

\textbf{Результат контроля.} Перемешанного или случайного fMRI в обучении нет. Случайные веса — только нуль для оценки предсказания мозга (не контрольная замена сигнала A)

\textbf{Другие сравнения.} исходная SoundNet без дообучения; индивидуальные против групповых моделей (HEAR: разницы нет); модель одного испытуемого на другом (только предсказание мозга, не HEAR); 29 чужих моделей из таблицы HEAR

\textbf{Объём чтения.} полный текст: препринт bioRxiv 2023 (arXiv-аналог, PDF) и версия PMC/EuropePMC XML (Imaging Neuroscience 2025); сайт MIT Press отдал заглушку Cloudflare, не использован


\Needspace{9\baselineskip}\subsection*{26. \href{https://doi.org/10.21437/interspeech.2025-2776}{Brain-tuned Speech Models Better Reflect Speech Processing Stages in the Brain}}

\textbf{Год и группа.} 2025; группа 2; пограничная: нет

\textbf{Обучающий сигнал A.} данные fMRI (Moth Radio Hour, 25 рассказов, 8 человек): дообучение речевых моделей предсказывать измеренную активность мозга

\textbf{Оценка B.} послойное выравнивание с поздними языковыми областями мозга и пробы на независимых наборах (MFCC, слова, фонемы TIMIT, фонетическая категория предложения TIMIT: SA/SX/SI)

\textbf{Модель и данные.} Wav2Vec2.0 и HuBERT base; по одной модели на участника

\textbf{Утверждение.} Aligning model representations more closely with brain data not only boosts performance on high-level downstream tasks

\textbf{Место в тексте.} §4, с. 4 PDF (заключение); аннотация, с. 1

\textbf{Контроль.} не найден

\textbf{Результат контроля.} Перемешанных или случайных данных мозга в обучении нет

\textbf{Другие сравнения.} исходные предобученные модели (основное сравнение)

\textbf{Объём чтения.} полный текст (аннотация, §1–§4)


\Needspace{9\baselineskip}\subsection*{27. \href{https://doi.org/10.48550/arxiv.2202.13100}{SemSup: Semantic Supervision for Simple and Scalable Zero-shot Generalization}}

\textbf{Год и группа.} 2022; группа 1; пограничная: да

\textbf{Обучающий сигнал A.} описания классов (несколько на класс, текст или JSON) как семантический надзор; обучение на видимых классах

\textbf{Оценка B.} классы и суперклассы, не использованные при обучении (ZSL) на 4 наборах: текст и изображения

\textbf{Модель и данные.} SemSup (CLIP/T5, гибридное сходство); CIFAR100, AwA2 и др.

\textbf{Утверждение.} One way to achieve this generalization is integrating class semantics into the learning framework

\textbf{Место в тексте.} §1, с. 1 PDF; аннотация: SemSup increases unseen class generalization accuracy by 15 points on average compared to the closest baseline; табл. 6, с. 7

\textbf{Контроль.} не найден

\textbf{Результат контроля.} Пограничный случай (косвенная связь результата с содержанием A: улучшение приписано описаниям, а не изучено, что именно в них важно). Контрольных описаний (перемешанных между классами или случайных) нет

\textbf{Другие сравнения.} абляции: только текст аннотаций, только изображения, одно против нескольких описаний, JSON против текста; CLIP как базовая модель; слой исправлен с 2 на 1 (по образцу InfZSL)

\textbf{Объём чтения.} полный текст (аннотация, §1, §4, §5, §7, табл. 6)


\Needspace{9\baselineskip}\subsection*{28. \href{https://doi.org/10.18653/v1/2021.cmcl-1.26}{Relation Classification with Cognitive Attention Supervision}}

\textbf{Год и группа.} 2021; группа 2; пограничная: нет

\textbf{Обучающий сигнал A.} данные взгляда (ET) и ЭЭГ (EEG) корпуса ZuCo как цель для внимания BERT при дообучении классификатору отношений

\textbf{Оценка B.} внимание на отложенных 57 примерах (пересечение топ-k слов с человеческим вниманием, KL), уникальные ошибки и исправления относительно базовой модели

\textbf{Модель и данные.} BERT; ZuCo, классификация отношений; 4 прогона

\textbf{Утверждение.} we find that this attention supervision can be used to increase similarity between model attention distributions over sequences and the cognitive data

\textbf{Место в тексте.} аннотация, с. 1 PDF; §5.2–5.3, табл. 4–5, с. 7–8; §6, с. 8

\textbf{Контроль.} есть

\textbf{Результат контроля.} Random ET / EEG / ET+EEG: значения ZuCo заменены равномерным шумом в том же диапазоне (§5.1, с. 6). Точность от надзора не меняется (табл. 3). Табл. 4: у ET настоящий надзор даёт больше пересечения с ET (26,1 против 6,8 у Random ET при k=1), у EEG контрольные метки не дают более низкого результата, чем настоящие (Random EEG 13,1 против EEG 9,5 при k=1 по метрике EEG). Табл. 5: настоящие данные чаще исправляют ошибки базовой модели, чем случайные. Вывод: замена обучающего сигнала проведена; для ET различие есть, для EEG — нет

\textbf{Другие сравнения.} базовая модель без надзора вниманием (дообучена на ZuCo); отдельные ET, EEG, ET+EEG

\textbf{Объём чтения.} полный текст (аннотация, §1, §3–§6, табл. 3–5)


\Needspace{9\baselineskip}\subsection*{29. \href{https://doi.org/10.1038/s42256-026-01278-w}{Beyond representational alignment with brain-guided language models for robust reasoning}}

\textbf{Год и группа.} 2026; группа 2; пограничная: нет

\textbf{Обучающий сигнал A.} task-fMRI 10 людей при дедуктивном рассуждении (силлогизмы, транзитивность); подсказка направления в представлениях LLM (NARI) и дообучение (NARF)

\textbf{Оценка B.} новые задачи: перестановка посылок, больше посылок, пропозициональное рассуждение, FOLIO; перенос общего направления на новые задачи

\textbf{Модель и данные.} 10 LLM (1,5–72 млрд); fMRI-набор Monti и др. [не сверено], псевдослова

\textbf{Утверждение.} Ablation analyses confirm that the gains arise from the joint influence of model and human neural representational structures

\textbf{Место в тексте.} §2.2.3, с. 9–10 PDF; прил. S11.2, рис. S8, с. 36–37; §2.2.1 рис. 4d, с. 8; аннотация, с. 2

\textbf{Контроль.} есть

\textbf{Результат контроля.} Random Signals: настоящий fMRI заменён случайными данными (NARI рис. 4d, с. 8; RSRF в NARF рис. S8, с. 36–37). NARI: случайные сигналы дают ограниченный эффект, настоящие — 100\% исправлений. NARF: RSRF+Label тоже улучшает чисто метки (из-за информации в представлениях модели), но настоящий fMRI значимо больше. Вид «случайных данных» (распределение) — [не сверено]; перемешанного fMRI между задачами нет

\textbf{Другие сравнения.} случайные направления; метки в виде one-hot (LSRF) — не лучше RSRF; только языковой надзор (метки); Rand. Sig. с средним fMRI; области мозга (рассуждение против языка и MD)

\textbf{Объём чтения.} полный текст (аннотация, §2.2, методы, прил. S5, S10.2, S11.2, рис. S5, S8)


\Needspace{9\baselineskip}\subsection*{30. \href{https://doi.org/10.52202/079017-0952}{Learning Human-like Representations to Enable Learning Human Values}}

\textbf{Год и группа.} 2024; группа 2; пограничная: да

\textbf{Обучающий сигнал A.} суждения людей о сходстве (человеческие и от языковых моделей) как ядро представлений агента; граница: ядро не обучается, а задаётся

\textbf{Оценка B.} оценка нравственности и ещё 10 ценностей на 25 действиях, не использованных при обучении (фаза обобщения), безопасность исследования

\textbf{Модель и данные.} многорукий бандит с гауссовым процессом; текстовые описания действий; суждения 20 людей

\textbf{Утверждение.} demonstrate that representational alignment enables both safe exploration and improved generalization when learning human values

\textbf{Место в тексте.} аннотация, с. 1 PDF; §6, с. 10 (strong evidence); контроль §5, с. 10 и прил. A.9, с. 22; §4.1, с. 6–7

\textbf{Контроль.} частичный

\textbf{Результат контроля.} Пограничный случай (A не обучает модель, а задаёт ядро). Моделирование: у k действий оценка нравственности заменена случайной, ядро портится, выравнивание падает, результат ухудшается (§4.1, с. 7). Контроль с другим осмысленным сигналом: ядро по длине описания — человеческое ядро сильно лучше при нравственности (прил. A.9, с. 22). Перемешанных суждений людей между действиями нет

\textbf{Другие сравнения.} разное число модельных ядер (языковые модели с разным выравниванием); контроль с другой функцией награды (длина); сложность задачи — без значимого эффекта

\textbf{Объём чтения.} полный текст (аннотация, §1–§6, прил. A.8–A.9)


\Needspace{9\baselineskip}\subsection*{31. \href{https://doi.org/10.48550/arxiv.2405.19420}{Learning Human-Aligned Representations with Contrastive Learning and Generative Similarity}}

\textbf{Год и группа.} 2024; группа 2; пограничная: да

\textbf{Обучающий сигнал A.} порождающее сходство (Байес, вероятностные программы) в контрастной потере; не человеческие данные, а теоретическая замена суждений людей

\textbf{Оценка B.} человеческие результаты: ошибки людей и бабуинов в задаче «лишнее», ряд геометрических понятий (7 отложенных из 37), суждения людей Levels-набора на естественных изображениях

\textbf{Модель и данные.} CorNet, ResNet; квадрилатерали, Geoclidean, ImageNet

\textbf{Утверждение.} We demonstrate the utility of our approach by showing that it can be used to capture human-like representations of shape regularity

\textbf{Место в тексте.} аннотация, с. 1 PDF; рис. 2–4; §3.1, с. 5; прил. O, с. 21 (30 понятий на обучение, 7 на тест)

\textbf{Контроль.} не найден

\textbf{Результат контроля.} Пограничный случай (A — синтетический сигнал сходства, не данные людей). Бессмысленного сигнала (случайные сходства / перемешанные программы) нет

\textbf{Другие сравнения.} SimCLR (без знания домена); потеря по классам (SupCon-подобная); исходная предобученная модель

\textbf{Объём чтения.} полный текст (аннотация, §1–§4, прил. L, O, P)


\Needspace{9\baselineskip}\subsection*{32. \href{https://doi.org/10.48550/arxiv.2510.01502}{Behavioral Geometric Supervision Aligns Video Foundation Models with Human Social Perception}}

\textbf{Год и группа.} 2025; группа 2; пограничная: да

\textbf{Обучающий сигнал A.} 49 484 суждения людей «лишнее» по 250 видео социального взаимодействия (потеря BGS: тройки + RSA на сходство)

\textbf{Оценка B.} пять свойств социальных сцен (общение, эмоциональная валентность, возбуждение, близость и доминирование), оцениваемых линейными пробами по фиксированным видеоэмбеддингам; их метки не использовали при дообучении сети, нулевой перенос на абстрактные взаимодействия PHASE, сохранение распознавания действий (UCF101)

\textbf{Модель и данные.} V-JEPA 2/2.1, TimeSformer, VideoMAE, CLIP; LoRA

\textbf{Утверждение.} Because they never appear in the supervision signal, such gains indicate that human similarity judgments encode an interpretable social-affective geometry that BGS makes linearly accessible

\textbf{Место в тексте.} §4.4, с. 7 PDF; аннотация, с. 1; §4.3, с. 6 (контроль); заключение, с. 9

\textbf{Контроль.} частичный

\textbf{Результат контроля.} Пограничный случай. Контроль «языковая дистилляция»: та же функция потерь и те же веса, но цели и тройки взяты из сходства подписей MPNet вместо суждений людей (§3, с. 5). Он смысловой, не случайный. Результат: R2 0,123 против 0,172 у BGS; атрибуты — контроль не воспроизводит прирост на 4 из 5. Перемешанных суждений людей или случайных троек в обучении нет (permutation на оценке — другое)

\textbf{Другие сравнения.} исходные предобученные модели; только тройки / только RSA; тройки того же бюджета; MPNet как внешний потолок; разные основы

\textbf{Объём чтения.} полный текст (аннотация, §1, §3–§4, §6, прил. H–I)


\Needspace{9\baselineskip}\subsection*{33. \href{https://doi.org/10.63317/5f2if588t6s5}{Linguistic Knowledge Can Enhance Encoder-Decoder Models (If You Let It)}}

\textbf{Год и группа.} 2024; группа 2; пограничная: да

\textbf{Обучающий сигнал A.} структурные лингвистические свойства предложений (10 признаков, предсказание как промежуточная задача, мультизадачно или по одному)

\textbf{Оценка B.} сложность предложения (оценка людей, корреляция Спирмена) на итальянском и английском, в т. ч. межъязыково

\textbf{Модель и данные.} T5 / IT5 / mT5 разных размеров; итальянские и английские наборы сложности предложений

\textbf{Утверждение.} linguistically motivated intermediate fine-tuning has generally a positive impact on target task performance

\textbf{Место в тексте.} аннотация, с. 1 PDF; заключение, с. 13-14 (§5): улучшение приписано «linguistic knowledge». Пограничный случай: косвенная связь результата с содержанием A

\textbf{Контроль.} не найден

\textbf{Результат контроля.} Случайных/перемешанных значений признаков в промежуточной задаче нет (поиск random/shuffl по тексту — пусто)

\textbf{Другие сравнения.} та же T5 без промежуточной стадии; разные размеры, число эпох, доля данных, межъязыковая схема

\textbf{Объём чтения.} полный текст (аннотация, введение, §3-5, ограничения)


\Needspace{9\baselineskip}\subsection*{34. \href{https://doi.org/10.17605/osf.io/pwr3t}{Intermediate -Task Transfer Learning: Leveraging Sarcasm Detection for Stance Detection}}

\textbf{Год и группа.} 2024; группа 2; пограничная: нет

\textbf{Обучающий сигнал A.} обнаружение сарказма (промежуточная задача; наборы ST, Sav2C, SARC)

\textbf{Оценка B.} определение позиции (stance detection): SemEval 2016 Task 6A и MPCHI

\textbf{Модель и данные.} BERT/RoBERTa + Conv + BiLSTM; arXiv 2503.03172 (найден по названию; DOI из входа — OSF)

\textbf{Утверждение.} the infusion of sarcasm knowledge significantly enhances model performance on the SemEval task

\textbf{Место в тексте.} §Results, G.4 «Ablation study regarding sarcasm knowledge», с. 5 PDF; табл. III-IV; аннотация, с. 1

\textbf{Контроль.} не найден

\textbf{Результат контроля.} Пустышки (сарказм со случайными метками) нет. Показано, что не всякая промежуточная задача помогает (Sav2C, SARC ухудшают), объяснено несовпадением данных

\textbf{Другие сравнения.} модель без сарказм-предобучения (табл. IV, BERT+Conv+BiLSTM 0.725/0.721 против 0.775/0.724); абляция слоёв; сравнение с SOTA

\textbf{Объём чтения.} полный текст


\Needspace{9\baselineskip}\subsection*{35. \href{https://doi.org/10.48550/arxiv.2608.27877}{Relational Knowledge Distillation Brings DNN Representations Close Enough to Humans to Be Aligned Without Supervision}}

\textbf{Год и группа.} 2026; группа 2; пограничная: нет

\textbf{Обучающий сигнал A.} человеческое психологическое вложение по суждениям о сходстве (THINGS), перенос парных отношений (RKD)

\textbf{Оценка B.} выравнивание на отложенном тестовом наборе из 300 понятий, не пересекающихся с обучением, без учителя (Gromov-Wasserstein); глобальная структура по категориям

\textbf{Модель и данные.} CLIP ViT-B/16 и ещё 4 варианта; DINO, supervised, ViT со случайной инициализацией; ImageNet-образцы

\textbf{Утверждение.} relational transfer from humans brings the global structure of pre-trained DNNs close enough to the human structure to enable fine-grained human–DNN alignment

\textbf{Место в тексте.} аннотация, с. 1 PDF (улучшение «driven by a more human-like global structure»); §3.1, прил. S3

\textbf{Контроль.} не найден

\textbf{Результат контроля.} Перемешанного / случайного учителя нет (поиск shuffl — пусто)

\textbf{Другие сравнения.} модель до RKD; случайно инициализированный ViT (после RKD остаётся около шанса); уровни шанса; контроль «деление внутри THINGS» (прил. S3)

\textbf{Объём чтения.} полный текст (аннотация, введение, методы, §3, прил.)


\Needspace{9\baselineskip}\subsection*{36. \href{https://doi.org/10.48550/arxiv.2604.20027}{Cognitive Alignment At No Cost: Inducing Human Attention Biases For Interpretable Vision Transformers}}

\textbf{Год и группа.} 2026; группа 2; пограничная: нет

\textbf{Обучающий сигнал A.} человеческие карты фиксаций взгляда (дообучение весов самовнимания)

\textbf{Оценка B.} не обучавшиеся свойства: предпочтение мелких объектов, одушевлённости, энтропия внимания; выравнивание на отложенных изображениях; классификация ImageNet/-C/ObjectNet

\textbf{Модель и данные.} ViT-B/16 (Google); повтор на ResNet-50

\textbf{Утверждение.} this confirms that the alignment gain reflects the structured, semantic content of the fixation signal rather than a generic effect of additional human supervision

\textbf{Место в тексте.} §Results, «Saliency Alignment», с. 4 PDF; рис. 2-3; обсуждение, с. 11

\textbf{Контроль.} есть

\textbf{Результат контроля.} Пары «изображение — карта фиксаций» перемешаны (seed 42), объём разметки и функция потерь сохранены. Обучение на настоящих картах превосходит контроль по всем 5 метрикам выравнивания. Изменение направления предпочтения мелких объектов найдено только при настоящих картах. Для предпочтения одушевлённых объектов: исходная модель Δ=+0.021, контроль +0.025, настоящие карты +0.029. Изменение энтропии: контроль -0.012, настоящие карты -0.043 бит. Часть изменений наблюдается и при обучении на перемешанных картах.

\textbf{Другие сравнения.} исходная ViT; ResNet-50 (выравнивание и точность ухудшились); байесовская проверка равенства точности

\textbf{Объём чтения.} полный текст


\Needspace{9\baselineskip}\subsection*{37. \href{https://doi.org/10.48550/arxiv.2505.14204}{Beginning with You: Perceptual-Initialization Improves Vision-Language Representation and Alignment}}

\textbf{Год и группа.} 2025; группа 2; пограничная: нет

\textbf{Обучающий сигнал A.} человеческие суждения о сходстве по тройкам (NIGHTS, \textasciitilde{}20 тыс. троек), инициализация кодировщика изображений

\textbf{Оценка B.} zero-shot классификация на 29 наборах и поиск на 2 (MS-COCO, Flickr30k), не входившие в обучение

\textbf{Модель и данные.} CLIP ViT-B/32; затем самообучение на YFCC15M

\textbf{Утверждение.} seeding models with human perceptual priors fosters systematically enhanced generalization capabilities on a wide spectrum of unseen classification tasks

\textbf{Место в тексте.} §4.1, с. 6 PDF; §4.4 «indicating that the human-derived triplet supervision indeed steers the representation toward more semantically faithful matches», с. 8; аннотация

\textbf{Контроль.} не найден

\textbf{Результат контроля.} Перемешанных или случайных троек нет. Контроль — модель со случайной инициализацией (без A)

\textbf{Другие сравнения.} та же модель с нуля на YFCC15M; «human-later» (то же дообучение после предобучения) — ломает выравнивание; планируемое, но не сделанное исследование доли троек

\textbf{Объём чтения.} полный текст


\Needspace{9\baselineskip}\subsection*{38. \href{https://doi.org/10.1016/j.neuron.2018.03.044}{A Task-Optimized Neural Network Replicates Human Auditory Behavior, Predicts Brain Responses, and Reveals a Cortical Processing Hierarchy}}

\textbf{Год и группа.} 2018; группа 2; пограничная: нет

\textbf{Обучающий сигнал A.} метки задач распознавания слов (речь) и жанра (музыка) по звуку

\textbf{Оценка B.} отклики мозга (fMRI вокселей слуховой коры) на естественные звуки и ошибки человека-слушателя; в обучении не участвовали

\textbf{Модель и данные.} двухветвевая свёрточная сеть по кохлеограмме; звуки с метками слов и жанров; fMRI людей

\textbf{Утверждение.} exhibited human-like errors despite not being optimized to do so, suggesting common constraints on network and human performance

\textbf{Место в тексте.} аннотация, с. 1 PDF (рис. 2 — поведение, рис. 3-4 — fMRI)

\textbf{Контроль.} не найден

\textbf{Результат контроля.} Результат не зафиксирован в исходной таблице извлечения.

\textbf{Другие сравнения.} сеть с теми же архитектурой и случайными необученными весами (рис. 3B: предсказывает fMRI заметно хуже обученной); спектрально-временная модель корковых фильтров; перемешанных меток слов/жанров нет

\textbf{Объём чтения.} полный текст (аннотация, введение, результаты, подписи к рисункам)


\Needspace{9\baselineskip}\subsection*{39. \href{https://doi.org/10.1038/s41467-023-43713-1}{Augmenting interpretable models with large language models during training}}

\textbf{Год и группа.} 2023; группа 2; пограничная: да

\textbf{Обучающий сигнал A.} бинарные метки тональности предложений (SST2 и др.); знания LLM входят через эмбеддинги n-грамм (BERT/LLaMA)

\textbf{Оценка B.} коэффициенты для не виденных в обучении триграмм; оценки тональности фраз людьми (SST), не использовавшиеся в обучении

\textbf{Модель и данные.} Aug-GAM (обобщённая аддитивная модель на эмбеддингах LLM), Aug-Tree; SST2, Rotten Tomatoes, Financial Phrasebank, Emotion

\textbf{Утверждение.} The inferred coefficients are semantically meaningful, even capturing three-way interactions, such as not very amusing

\textbf{Место в тексте.} §3.4 (раздел «Inferred Aug-GAM coefficients for unseen ngrams match human scores»), с. 8 PDF; рис. 4C-D

\textbf{Контроль.} не найден

\textbf{Результат контроля.} Единственное «случайное спаривание» — базовая оценка людей для расширений n-грамм Aug-Tree (табл. 3: 1.3 против 4.6), это контроль оценки, а не обучающего сигнала. Пограничный случай: A — знание LLM, а не обучение на человеческом сигнале; включена по критерию косвенной связи результата с содержанием A

\textbf{Другие сравнения.} мешок биграмм без LLM (ρ 0.39 против 0.63; 0.49 против 0.71 на триграммах); модели без аугментации; большие LLM (GPT-J)

\textbf{Объём чтения.} полный текст (аннотация, §2-3, табл. 3, рис. 4); версия arXiv v3


\Needspace{9\baselineskip}\subsection*{40. \href{https://doi.org/10.48550/arxiv.2310.06089}{Predictive auxiliary objectives in deep RL mimic learning in the brain}}

\textbf{Год и группа.} 2023; группа 2; пограничная: да

\textbf{Обучающий сигнал A.} вспомогательная предсказательная потеря (предсказание следующего латентного состояния, контрастная L+ и L-) при обучении агента RL в задаче A

\textbf{Оценка B.} перенос на задачу B (новая цель, перемешанные переходы среды) с замороженным кодировщиком; сходство с данными мозга из литературы (зрительная кора, гиппокамп, полосатое тело)

\textbf{Модель и данные.} агент глубокого RL (кодировщик + обучение ценности + модуль предсказания); игрушечные лабиринты 8x8, круговая дорожка, T-образный лабиринт

\textbf{Утверждение.} We identify settings where longer predictive horizons better support representational transfer

\textbf{Место в тексте.} аннотация, с. 1 PDF; рис. 3, с. 6

\textbf{Контроль.} не найден

\textbf{Результат контроля.} Есть «случайное перемешивание» (рис. 4G, с. 7) — контроль для сопоставления рецептивных полей с мозгом (перемешаны веса модуля T), а не обучающего сигнала

\textbf{Другие сравнения.} модель без вспомогательных потерь; вариант только с L- (рис. 3G-H); разные горизонты γ (γ = 0); разные размеры латентного слоя. Пограничный случай: приписывание — «предсказательной цели», не смыслу меток; включена по критерию косвенной связи результата с содержанием A

\textbf{Объём чтения.} аннотация, введение, §3-4 (рис. 3-4), подписи; приложение не читалось


\Needspace{9\baselineskip}\subsection*{41. \href{https://doi.org/10.18653/v1/2022.findings-naacl.156}{Beyond Distributional Hypothesis: Let Language Models Learn Meaning-Text Correspondence}}

\textbf{Год и группа.} 2022; группа 2; пограничная: нет

\textbf{Обучающий сигнал A.} промежуточная задача «meaning-matching»: соответствует ли определение слова (WordNet-подобные пары) самому слову; ложные пары получены негативной выборкой

\textbf{Оценка B.} зонды лексической семантики и отрицания (SAR, MKR-NQ, MWR), NegNLI и 7 задач GLUE — на них модель не обучалась

\textbf{Модель и данные.} BERT/RoBERTa/ELECTRA (small–large); дообучение на meaning-matching, затем зонды

\textbf{Утверждение.} we find that the task enables PLMs to learn lexical semantic information

\textbf{Место в тексте.} аннотация, с. 1 PDF; табл. 4-5, с. 6; табл. 7, с. 7

\textbf{Контроль.} не найден

\textbf{Результат контроля.} Ложные пары определений — часть самой задачи (отрицательные примеры), не контрольная замена правильных пар; перемешанных определений вместо настоящих нет

\textbf{Другие сравнения.} исходные PLM до промежуточного обучения (табл. 4, 5, 7); BERTNOT (табл. 9); вариант с замороженным кодировщиком; оценка забывания (норма Фробениуса)

\textbf{Объём чтения.} полный текст (аннотация, введение, §2-5, табл. 3-9)


\Needspace{9\baselineskip}\subsection*{42. \href{https://doi.org/10.48550/arxiv.2410.02915}{Does the Order of Fine-tuning Matter and Why?}}

\textbf{Год и группа.} 2024; группа 2; пограничная: да

\textbf{Обучающий сигнал A.} промежуточные задачи дообучения CodeBERT (клон-детекция, обнаружение дефектов, исправление и перевод кода), в разных порядках цепочек

\textbf{Оценка B.} целевая задача в конце цепочки (качество на ней), а также зонды семантики и синтаксиса кода

\textbf{Модель и данные.} CodeBERT; задачи CodeXGLUE (CD, DD, CR, CT)

\textbf{Утверждение.} Choosing datasets that are more semantically similar could lead to better performance

\textbf{Место в тексте.} §5.2, с. 12 PDF (наблюдение после анализа семантической близости наборов)

\textbf{Контроль.} не найден

\textbf{Результат контроля.} Результат не зафиксирован в исходной таблице извлечения.

\textbf{Другие сравнения.} модели, дообученные только на целевой задаче; контроль размера наборов (FULL/LIMITED, случайная подвыборка 1K, табл. 3-4). Пограничный случай: приписывание — близости наборов A и B, а не смыслу A

\textbf{Объём чтения.} полный текст (аннотация, введение, §4-5, табл. 3-4)


\Needspace{9\baselineskip}\subsection*{43. \href{https://doi.org/10.21203/rs.3.rs-11089803/v1}{Extremely coarse learning objectives induce human-aligned representations in AI vision models}}

\textbf{Год и группа.} 2026; группа 2; пограничная: нет

\textbf{Обучающий сигнал A.} метки грубых классов (2–64 класса), полученные делением обучающих изображений по главным компонентам эмбеддингов (CLIP, AlexNet и др.); также 8 вручную заданных семантических классов (естественное/искусственное, размер, помещение/улица), классы WordNet

\textbf{Оценка B.} выравнивание с fMRI человека (NSD), электрофизиологией макаки и суждениями людей о сходстве (THINGS, тройки «лишний»); в обучении не участвовали

\textbf{Модель и данные.} AlexNet-подобная сеть, ResNet-50, ConvNeXt-Base, ViT-B/16, обучены с нуля на ImageNet; 3 случайные инициализации

\textbf{Утверждение.} human-like visual representations can emerge from surprisingly simple learning objectives

\textbf{Место в тексте.} аннотация, с. 1 PDF; рис. 2, 3 и 6, с. 6-10 (arXiv 2605.05556v2)

\textbf{Контроль.} частичный

\textbf{Результат контроля.} Контроль «pixel PCA» (с. 6, методы с. 17): тоже 3 бита/K классов и та же функция потерь, но классы по низкоуровневым пикселям (яркость, контраст) — слабое выравнивание с нейронными данными. Это контроль на основе низкоуровневых свойств изображения при сходном формате классов; метки не перемешиваются

\textbf{Другие сравнения.} необученная сеть той же архитектуры (эпоха 0); 1000-классная сеть с теми же данными и рецептом; классы из разных исходных моделей; 8 ручных семантических классов; перемешанных/случайных меток нет. Важно: итоги по сходству с людьми почти не зависят от точности классификации

\textbf{Объём чтения.} полный текст (аннотация, введение, §2-6, 10, методы, подписи)


\Needspace{9\baselineskip}\subsection*{44. \href{https://doi.org/10.1007/978-3-030-61609-0_69}{Recurrent Neural Network Learning of Performance and Intrinsic Population Dynamics from Sparse Neural Data}}

\textbf{Год и группа.} 2020; группа 2; пограничная: да

\textbf{Обучающий сигнал A.} разреженные записи активности подмножества нейронов биологической (модельной) сети плюс выходные сигналы (ЭМГ)

\textbf{Оценка B.} обобщение на отложенные направления движения (ошибка обобщения) и сходство внутренней динамики популяции (канонические корреляции)

\textbf{Модель и данные.} RNN, обучение Full-FORCE; физиологически мотивированная модель моторной коры (центр–наружу), синтетические данные

\textbf{Утверждение.} training the RNNs with this method significantly improves their generalization performance

\textbf{Место в тексте.} аннотация, с. 1 PDF; §3 «Embedding the Dynamics Improves RNN ability to Generalize», с. 8; рис. 5

\textbf{Контроль.} не найден

\textbf{Результат контроля.} Базовая линия CC — необученная RNN со случайными связями (D0), но это не контрольная замена обучающего сигнала

\textbf{Другие сравнения.} RNN, обученные только на выходе (без внедрения динамики); необученная случайная RNN. Пограничный случай: данные синтетические, приписывание — «захватывает вычислительные свойства» схемы

\textbf{Объём чтения.} аннотация, введение, §2.4, §3, выводы (препринт arXiv v1)


\Needspace{9\baselineskip}\subsection*{45. \href{https://doi.org/10.48550/arxiv.2610.22569}{Alignment of visual brain encoding models with individual perceptual similarity judgments}}

\textbf{Год и группа.} 2026; группа 2; пограничная: нет

\textbf{Обучающий сигнал A.} суждения людей о сходстве 100 изображений (индивидуальные, 8 участников NSD) $\rightarrow$ персонализированный DreamSim $\rightarrow$ потеря RSA при дообучении кодировщика мозга на \textasciitilde{}9000 обучающих изображений

\textbf{Оценка B.} согласие кодировщиков с измеренными суждениями участника о сходстве (отложенные пары, 5-кратная перекрёстная проверка, 10 повторов); межсубъектная структура

\textbf{Модель и данные.} NRF-кодировщики мозга (fMRI), NSD; DreamSim; 8 участников

\textbf{Утверждение.} fine-tuning brain encoding models with personalized perceptual similarity information results in models that are overall the most aligned with human perceptual similarity judgments

\textbf{Место в тексте.} аннотация, с. 1 PDF; §5, с. 9 (transferring these properties to the personalized brain encoders)

\textbf{Контроль.} не найден

\textbf{Результат контроля.} Результат не зафиксирован в исходной таблице извлечения.

\textbf{Другие сравнения.} исходный кодировщик; тот же кодировщик с общими (популяционными) суждениями DreamSim при той же инициализации и потере; перестановочный тест подписей участников — только для статистики; перемешанных или чужих суждений как обучающего сигнала нет

\textbf{Объём чтения.} полный текст (аннотация, введение, §3, §4, §5)


\Needspace{9\baselineskip}\subsection*{46. \href{https://doi.org/10.48550/arxiv.2606.28399}{Meta-learning as a principle for human-like visual representations}}

\textbf{Год и группа.} 2026; группа 2; пограничная: да

\textbf{Обучающий сигнал A.} мета-обучение на тысячах задач «изображение $\rightarrow$ активация латента SAE (понятие)»; обучающих данных людей нет

\textbf{Оценка B.} предсказание суждений людей о сходстве, обучение семантическим правилам, отклик зрительной коры (5 открытых наборов)

\textbf{Модель и данные.} каузальный трансформер + линейная проекция признаков SigLIP2, MAE, DINOv3, CLIP; THINGS, NSD и др.

\textbf{Утверждение.} Our results suggest the flexibility of human visual representations reflects the functional demand to learn new semantic relationships on the fly.

\textbf{Место в тексте.} аннотация, с. 1 PDF; «The role of disentanglement», с. 6, рис. 3–4

\textbf{Контроль.} частичный

\textbf{Результат контроля.} Контроль «сырой остаточный поток CLIP» (запутанные задачи того же уровня) и задачи среднего уровня (слой 6) дают хуже на поведении; на мозге картина неоднородна (высокоуровневое не систематически лучше). Случайных или перемешанных задач нет — это сравнение альтернативных распределений задач

\textbf{Другие сравнения.} исходные замороженные кодировщики; многозадачная модель на тех же задачах без эпизодической структуры (рис. 3)

\textbf{Объём чтения.} полный текст (аннотация, введение, результаты, методы, дополнение)


\Needspace{9\baselineskip}\subsection*{47. \href{https://doi.org/10.48550/arxiv.2609.36239}{Cognitive Expert Language Models Better Align with the Corresponding Brain Systems}}

\textbf{Год и группа.} 2026; группа 2; пограничная: нет

\textbf{Обучающий сигнал A.} дообучение (LoRA) или инструкция-подсказка на тексте одного из 6 когнитивных доменов (сенсорный, пространственный, числовой, рассуждение, социальный, абстрактный)

\textbf{Оценка B.} согласие представлений модели с активностью отдельных областей мозга (fMRI, 3 набора): «эксперт» лучше предсказывает свою область

\textbf{Модель и данные.} LLaMA-3-8B, Qwen3-8B, Phi-4-14B; три набора fMRI (в т. ч. Le Petit Prince, LeBel, Pereira)

\textbf{Утверждение.} each expert’s representations align more closely with the brain system most associated with the matching domain than do other experts

\textbf{Место в тексте.} аннотация, с. 1 PDF; §3.1 «Control models», с. 4; §4.4, с. 8, рис. 5; прил. C

\textbf{Контроль.} есть

\textbf{Результат контроля.} Пять семейств-контролей с той же процедурой без когнитивного содержания: случайные LoRA, 6 адаптеров на одном общем смешанном корпусе (различаются только зерном), случайные подсказки («анатомия полосатости зебры»), адаптеры на тексте с искажённой формой (leetspeak, пиратский английский), ключевые слова без инструкции. Эксперты в среднем на 94-м процентиле нулевого распределения «seed-only», 14 из 18 выше 95-го; у адаптеров с искажённой формой изменения в \textasciitilde{}4 раза больше, но регионального соответствия нет; ключевые слова слабее в 7 из 9 сравнений (исключения — Qwen на двух нарративных наборах)

\textbf{Другие сравнения.} базовая модель; сравнение подсказки и дообучения; 720 перестановок назначения «эксперт $\rightarrow$ область»

\textbf{Объём чтения.} полный текст (аннотация, §3–4, прил. C–D, H)


\Needspace{9\baselineskip}\subsection*{48. \href{https://doi.org/10.48550/arxiv.2505.23933}{BIRD: Behavior Induction via Representation-structure Distillation}}

\textbf{Год и группа.} 2025; группа 2; пограничная: да

\textbf{Обучающий сигнал A.} структура (матрицы Грама) промежуточных представлений учителя, обученного устойчивости (ImageNet-C), как сигнал для ученика на другом наборе

\textbf{Оценка B.} робастная точность ученика на искажённых данных (в обучении ученика только чистые)

\textbf{Модель и данные.} BIRD; CIFAR-10/100, TinyImageNet; AlexNet, MobileNetV2 и др. учителя, ResNet50 ученик; GPT-2 (§6)

\textbf{Утверждение.} a flexible framework for transferring aligned behavior by matching the internal representation structure of a student model to that of a teacher

\textbf{Место в тексте.} аннотация, с. 1 PDF; §5, с. 7–8, рис. 5 («behavioral relevance» объясняет >50\% дисперсии)

\textbf{Контроль.} не найден

\textbf{Результат контроля.} Результат не зафиксирован в исходной таблице извлечения.

\textbf{Другие сравнения.} сравнение с дообучением, Hints, непрерывным обучением; пул из 144 учителей, у которых содержание варьируют (набор данных, доля искажений) — график «доза–отклик», контрольная замена обучающего сигнала не проводится; случайного/необученного учителя не найдено [не сверено по всему тексту]

\textbf{Объём чтения.} полный текст (аннотация, введение, §3–6, часть приложений)


\Needspace{9\baselineskip}\subsection*{49. \href{https://doi.org/10.1038/s41586-023-06031-6}{Learnable latent embeddings for joint behavioural and neural analysis}}

\textbf{Год и группа.} 2023; группа 2; пограничная: да

\textbf{Обучающий сигнал A.} метки «поведения» (положение, направление; для видео — признаки DINO по кадрам) при контрастном обучении вложений нейронной активности

\textbf{Оценка B.} декодирование на отложенных повторах (положение; идентификатор кадра фильма в зрительной коре мыши)

\textbf{Модель и данные.} CEBRA; гиппокамп крысы, макака, Allen Brain Observatory (Neuropixels, 2-фотонная)

\textbf{Утверждение.} The DINO features themselves did not drive performance, as shuffling the features showed poor decoding

\textbf{Место в тексте.} «Decoding natural videos», с. 7 PDF, расширенный рис. S10j; перемешанные метки — с. 4, рис. 1–2, доп. рис. S5e

\textbf{Контроль.} есть

\textbf{Результат контроля.} Перемешанные метки (положение, DINO-признаки, перестановка по времени, то же распределение) дают слабое декодирование и более высокую потерю; перемешанные вложения — нуль для чисел Бетти

\textbf{Другие сравнения.} CEBRA-Time, pi-VAE, UMAP, tSNE, популяционный вектор + kNN

\textbf{Объём чтения.} полный текст препринта arXiv 2204.00673 (основной текст, расширенные данные)


\Needspace{9\baselineskip}\subsection*{50. \href{https://doi.org/10.1109/tpami.2023.3263181}{Decoding Visual Neural Representations by Multimodal Learning of Brain-Visual-Linguistic Features}}

\textbf{Год и группа.} 2023; группа 2; пограничная: нет

\textbf{Обучающий сигнал A.} фМРТ/ЭЭГ-ответы мозга на изображения seen-классов, совмещённые с визуальными признаками и текстами Википедии

\textbf{Оценка B.} декодирование классов, не использованных при обучении (novel) по мозговой активности (zero-shot, на этих классах мозговых данных в обучении нет)

\textbf{Модель и данные.} BraVL (мультимодальный генеративный VAE, mixture-of-product-of-experts + регуляризация по взаимной информации); три тримодальных набора (два Image-fMRI, один Image-EEG)

\textbf{Утверждение.} visual perception may be accompanied by linguistic influences to represent the semantics of visual stimuli

\textbf{Место в тексте.} аннотация, с. 1 PDF; заключение, с. 15 PDF (visual perception is accompanied by language influences)

\textbf{Контроль.} не найден

\textbf{Результат контроля.} Контроля бессмысленным текстом/перемешанными описаниями классов не найдено (поиск shuffl/permut/random/placebo/control)

\textbf{Другие сравнения.} сравнение: только визуальные признаки, только тексты, оба вместе; уровень случайного угадывания (chance); другие ZSL-методы

\textbf{Объём чтения.} полный текст (аннотация, введение, методы, опыты, заключение); arXiv 2210.06756


\Needspace{9\baselineskip}\subsection*{51. \href{https://doi.org/10.1371/journal.pcbi.1010932}{Robust deep learning object recognition models rely on low frequency information in natural images}}

\textbf{Год и группа.} 2023; группа 2; пограничная: да

\textbf{Обучающий сигнал A.} регуляризация сходства представлений по нейронным данным (V1 мыши, V1 макаки) при обучении классификатора (модели из прежних работ [7,8])

\textbf{Оценка B.} устойчивость к искажениям и состязательным атакам (не обучалась напрямую)

\textbf{Модель и данные.} ResNet (серые CIFAR10, мышь) и VGG19 (макака); сравнение с обычными моделями и 15 робастными моделями (RobustBench, ImageNet)

\textbf{Утверждение.} We hypothesize that the increased model robustness is partly due to the low spatial frequency preference inherited from the neural representation.

\textbf{Место в тексте.} аннотация, с. 1 PDF; вывод «strong evidence» — Results, Fig 3, с. 5 PDF. Пограничный случай: приписано свойству нейронного представления (низкие частоты), а не «смыслу»

\textbf{Контроль.} не найден

\textbf{Результат контроля.} Контрольной замены нейронных данных (перемешанных или случайных данных мозга) в этой работе не найдено; контроль — обычная модель без нейронной регуляризации. Перемешивание в тексте — только случайные перестановки гибридных изображений

\textbf{Другие сравнения.} обычная модель того же типа; другие робастные модели (состязательное обучение, аугментации); размытие как защита

\textbf{Объём чтения.} полный текст (аннотация, результаты, обсуждение); PLoS Comput Biol (PMC/PLOS PDF)


\Needspace{9\baselineskip}\subsection*{52. \href{https://doi.org/10.48550/arxiv.2607.18237}{The Many Senses of Visual Similarity: A Text-Prompted Image Perceptual Metric}}

\textbf{Год и группа.} 2026; группа 2; пограничная: да

\textbf{Обучающий сигнал A.} человеческие суждения о сходстве (около 1 млн суждений, 25 тыс. триплетов изображений, 260 тыс. комбинаций триплет-аспект, свободный текст аспекта)

\textbf{Оценка B.} 17 внешних алгоритмов зрения/графики (2AFC), другие бенчмарки (BAPPS, NIGHTS, GeneCIS, QARE-Bench)

\textbf{Модель и данные.} Qwen3-VL-Embedding 8B (дообучена в TPIPS); собственный набор триплетов + внешний набор

\textbf{Утверждение.} We demonstrate that TPIPS aligns more closely with human perception and generalizes reliably beyond the training distribution.

\textbf{Место в тексте.} аннотация, с. 1 PDF; §5 Generalization, с. 8 PDF; прил. B.1. Пограничный случай: приписано человеческим данным («human-aligned»), смысл как таковой не обсуждается

\textbf{Контроль.} не найден

\textbf{Результат контроля.} Перемешанных суждений / случайных триплетов нет

\textbf{Другие сравнения.} недообученный бэкбон и VLM-базовые (Gemini, GPT, InternVL и др.), LPIPS, DreamSim; человеческий консенсус; абляции: замороженный бэкбон, без условия аспекта, доля данных (табл. 7)

\textbf{Объём чтения.} полный текст (аннотация, введение, §3–5, обсуждение, прил. B.1, B.3); arXiv 2607.18237v2


\Needspace{9\baselineskip}\subsection*{53. \href{https://openalex.org/W7220748548}{Neural Data Needs Semantic Tokenization: Behavioral Events as Boundaries of Session-Transferable Tokens}}

\textbf{Год и группа.} 2026; группа 2; пограничная: нет

\textbf{Обучающий сигнал A.} поведенческие события испытания (начало стимула, начало движения) как границы токенов состояний; обучение декодированию поведения на мышах (IBL)

\textbf{Оценка B.} held-out сессии IBL (53), перенос замороженной модели на макак и Utah-массивы (направление движения — цель, которую границы не задают)

\textbf{Модель и данные.} TWS (tokenization with states) + CNN/Transformer/Mamba, сравнение с PerceiverIO/POYO; IBL, Steinmetz 2019, Gallego, NLB MC Maze

\textbf{Утверждение.} each state carries its own behavioral meaning

\textbf{Место в тексте.} аннотация, с. 1 PDF (название: «Neural Data Needs Semantic Tokenization»); контрольные замены — табл. 2, с. 7 PDF

\textbf{Контроль.} есть

\textbf{Результат контроля.} Контрольные границы: фиксированный бин, случайный бин, перемешанные между сессиями и внутри сессии, со сдвигом (jitter). Движение MCC: настоящие 0.565; перемешанные 0.199/0.196; случайные 0.171; фиксированные 0.215/0.188; сдвиг σ=5 даёт 0.523. Время события отдельной пробой без нейроданных: 0.908 при настоящих границах, 0.000 при контрольных границах — то есть само время события позволяет предсказывать целевую переменную; с контрольными границами TWS сохраняет около трети декодирования

\textbf{Другие сравнения.} TWS с случайной инициализацией (без обучения); PerceiverIO (токенизатор без границ); 6 бэкбонов; проба по одному времени события

\textbf{Объём чтения.} полный текст (аннотация, введение, §4.1–4.5, табл. 1–4); arXiv 2610.03001


\Needspace{9\baselineskip}\subsection*{54. \href{https://doi.org/10.64898/2026.08.19.745390}{LAND: Latent Aligned Neural-Behavioral Dynamics via Flow Matching for Generalizable Movement Decoding}}

\textbf{Год и группа.} 2026; группа 2; пограничная: да

\textbf{Обучающий сигнал A.} неспаренные кинематические траектории (поведенческие динамические приоры), выравнивание нейронных и поведенческих латентных динамик потоком

\textbf{Оценка B.} неизвестные скорости движения (zero-shot), другие конечности и субъекты (с малой дообучкой)

\textbf{Модель и данные.} LAND (flow matching); синтетические данные, эпидуральные BCI-записи участника с тетраплегией, МЕА-записи макак

\textbf{Утверждение.} These results support flow-based neural–behavioral alignment with unpaired kinematic priors as an approach for learning transferable neural representations and robust movement decoding

\textbf{Место в тексте.} аннотация, с. 1 PDF; абляции — табл. 5, с. 7 PDF. Пограничный случай: приписано поведенческому приору, смысл как таковой не обсуждается

\textbf{Контроль.} не найден

\textbf{Результат контроля.} Случайных/перемешанных траекторий нет. Есть «w/o auxiliary data» (без A), без flow-выравнивания, без факторного выравнивания, без разнообразия латентов. Приор «синтетические рандомизированные траектории» входит в данные, но как отдельный контроль не заявлен

\textbf{Другие сравнения.} сквозное обучение без вспомогательных данных; MSE-выравнивание вместо потока; POYO, NEDS как базовые методы

\textbf{Объём чтения.} полный текст (9 с.): аннотация, методы, результаты, абляции, ограничения; bioRxiv


\Needspace{24\baselineskip}\subsection*{55. \href{https://doi.org/10.48550/arxiv.2611.05231}{Diagnosis by Composition: Detecting and Localizing Unseen Diseases by Composing Learned Radiological Findings}}

\textbf{Год и группа.} 2026; группа 1; пограничная: да

\textbf{Обучающий сигнал A.} названия радиологических находок (findings), выровненные с признаками КТ; болезни как набор находок из руководств

\textbf{Оценка B.} шесть болезней, исключённых из обучения (leave-one-disease-out), и редкие болезни другого учреждения

\textbf{Модель и данные.} DbC (диагноз как фиксированная композиция карт находок); мультифазная КТ печени

\textbf{Утверждение.} the gain comes from composing findings rather than from the encoder or the architecture

\textbf{Место в тексте.} §4.2, с. 7 PDF; §4.3, с. 8 (интуиция: «Disease names may carry clinical meaning more than visual evidence»); аннотация, с. 1

\textbf{Контроль.} не найден

\textbf{Результат контроля.} Результат не зафиксирован в исходной таблице извлечения.

\textbf{Другие сравнения.} парные абляции Disease-VL (те же кодировщик, данные, потеря, но выравнивание с названиями болезней) и Disease-V (без языка); опубликованные SegVol, SAT, VoxTell. Бессмысленных/перемешанных находок нет. Пограничный случай: приписано структуре «находки вместо болезней», смысловая часть — интуиция без проверки

\textbf{Объём чтения.} полный текст (аннотация, введение, §3.3-4.3, обсуждение)
